\documentclass[12pt, a4paper]{article}

% ── Packages ──────────────────────────────────────────────────────────────────
\usepackage[utf8]{inputenc}
\usepackage[T1]{fontenc}
\usepackage{lmodern}
\usepackage{amsmath, amssymb}
\usepackage{mathtools}
\usepackage{booktabs}
\usepackage{array}
\usepackage{enumitem}
\usepackage{hyperref}
\usepackage[margin=2.5cm]{geometry}
\usepackage{parskip}
\usepackage{xcolor}
\usepackage{tcolorbox}
\usepackage{titlesec}
\usepackage{listings}

% ── Colours ───────────────────────────────────────────────────────────────────
\definecolor{intuitionbg}{RGB}{240,248,255}
\definecolor{simulationbg}{RGB}{255,250,240}
\definecolor{keybox}{RGB}{245,245,220}
\definecolor{historybg}{RGB}{245,240,250}
\definecolor{codebg}{RGB}{248,248,248}
\definecolor{codeframe}{RGB}{200,200,200}
\definecolor{codegreen}{RGB}{0,128,0}
\definecolor{codegray}{RGB}{128,128,128}

% ── Environments ──────────────────────────────────────────────────────────────
\newtcolorbox{intuition}{
  colback=intuitionbg, colframe=blue!40!black,
  title={\textbf{Intuitive Explanation}}, fonttitle=\bfseries
}
\newtcolorbox{simulation}{
  colback=simulationbg, colframe=orange!60!black,
  title={\textbf{Simulation}}, fonttitle=\bfseries
}
\newtcolorbox{keytakeaway}{
  colback=keybox, colframe=olive!60!black,
  title={\textbf{Key Takeaway}}, fonttitle=\bfseries
}
\newtcolorbox{plainlanguage}{
  colback=white, colframe=gray!50,
  left=6pt, right=6pt, top=4pt, bottom=4pt
}
\newtcolorbox{historynote}{
  colback=historybg, colframe=purple!40!black,
  title={\textbf{Historical Note}}, fonttitle=\bfseries
}

% ── Code listing style ───────────────────────────────────────────────────────
\lstdefinestyle{pythonstyle}{
  language=Python,
  basicstyle=\ttfamily\small,
  backgroundcolor=\color{codebg},
  frame=single,
  rulecolor=\color{codeframe},
  keywordstyle=\color{blue}\bfseries,
  commentstyle=\color{codegreen}\itshape,
  stringstyle=\color{red!70!black},
  numberstyle=\tiny\color{codegray},
  numbers=left,
  numbersep=6pt,
  breaklines=true,
  showstringspaces=false,
  tabsize=4,
  xleftmargin=1.5em,
  framexleftmargin=1.5em,
  aboveskip=1em,
  belowskip=1em,
  morekeywords={as, True, False, None},
}
\lstset{style=pythonstyle}

% ── Notation shorthand ────────────────────────────────────────────────────────
\newcommand{\E}{\mathbb{E}}
\newcommand{\Prob}{\mathbb{P}}
\newcommand{\Var}{\operatorname{Var}}
\newcommand{\Cov}{\operatorname{Cov}}
\newcommand{\Cor}{\operatorname{Cor}}
\newcommand{\given}{\,\vert\,}


% ══════════════════════════════════════════════════════════════════════════════
\title{\textbf{Sequential Testing for A/B Experiments}\\[6pt]
       \large A Practical Guide in Three Acts\\[4pt]
       \normalsize Version 2}
\author{}
\date{April 2026}

\begin{document}
\maketitle
\tableofcontents
\newpage

% ══════════════════════════════════════════════════════════════════════════════
\section*{Introduction}
\addcontentsline{toc}{section}{Introduction}

\textbf{The problem.}
Every company running A/B tests faces a temptation: peek at the results
before the experiment is over.  Traditional statistics breaks when you do
this --- the false positive rate inflates far beyond the promised 5\%.

\textbf{The solution.}
\emph{Sequential testing} provides statistical methods that remain valid
no matter when or how often you check.  Modern platforms like Eppo build
this into their pipeline automatically.

\textbf{This guide.}
Three acts, each with a simulation, intuitive explanations, and key formulas
with plain-language translations:
\begin{enumerate}[nosep]
  \item \textbf{The Peeking Problem} --- Why checking early is dangerous.
  \item \textbf{The Eppo Solution} --- How a modern platform solves it.
  \item \textbf{DIY Alternatives} --- What to do if you don't have Eppo.
\end{enumerate}

\textbf{Target audience.}
People comfortable with algebra and basic probability.  No calculus required.
Every formula is accompanied by a plain-language translation.

\textbf{What changed from v1:}
Removed subtitle from Act~3; added detailed implementation instructions
with step-by-step equations for each DIY method; added Appendix~A with
Python code for simulation and implementation snippets.


% ══════════════════════════════════════════════════════════════════════════════
% ACT 1
% ══════════════════════════════════════════════════════════════════════════════
\newpage
\section{Act~1 --- The Peeking Problem}

\subsection{The Motivation}

\begin{intuition}
You work at a company running an A/B test.  Half your users see the old
website, half see a new design.  Every day, you check the dashboard:
``Is the new design better?''  After a week, the dashboard shows $p=0.03$ ---
that looks significant!  You call it: the new design wins.

But here is the dirty secret of traditional statistics: \textbf{the more often
you check, the more likely you are to see a false positive.}  A test designed
to have a 5\% false positive rate can easily reach 20--30\% if you peek at it
repeatedly.  One in four or five ``winning'' experiments might actually be
noise.

This is not a theoretical concern.  It is the single most common statistical
mistake in online experimentation.
\end{intuition}

\subsection{The Simulation}

\begin{simulation}
\textbf{What is being simulated:}
1,000 A/B experiments where there is \emph{no real effect} --- any apparent
difference is pure noise.

Two groups of observations (control and treatment) are generated from
\emph{identical} distributions.  A standard $t$-test is computed after every
new observation arrives.

A counter shows: ``Number of times $p<0.05$ so far.''  The user watches the
$p$-value bounce around.  Every time it dips below $0.05$, it lights up red
--- a false positive.

\textbf{Controls:}
\begin{itemize}[nosep]
  \item How many total observations to simulate.
  \item How many independent experiments to run in parallel.
  \item Fixed seed toggle (deterministic replay).
\end{itemize}

A bar chart accumulates: ``Fraction of experiments that showed $p<0.05$ at
\emph{some} point during peeking.''  For 1,000 experiments with regular
peeking, this fraction climbs to 20--30\%.

A second mode checks only once, at the very end.  That fraction stays near
5\%.

\textbf{Key takeaway:}
Peeking inflates false positives.  The more you peek, the worse it gets.
\end{simulation}

\subsection{Why Does This Happen?}

\begin{intuition}
Traditional statistical tests promise: ``If there is no real effect, I will
falsely cry wolf at most 5\% of the time.''  But this promise comes with fine
print: \textbf{you may only look at the result once, at a pre-determined time.}

Every time you peek, you get another chance for randomness to fool you.  It is
like rolling a die repeatedly --- the more rolls, the more likely you are to
get a six eventually, even though any single roll has only a $1/6$ chance.
\end{intuition}

\subsection{How Bad Is It?}

\begin{center}
\renewcommand{\arraystretch}{1.3}
\begin{tabular}{@{} l c @{}}
\toprule
\textbf{Checking schedule} & \textbf{Actual false positive rate} \\
\midrule
Once at the end (as designed) & $\sim\!5\%$ \\
Daily for 1 week & $\sim\!13\%$ \\
Daily for 2 weeks & $\sim\!19\%$ \\
Daily for 4 weeks & $\sim\!25\%$ \\
Continuously (every observation) & Can exceed $30\%$ \\
\bottomrule
\end{tabular}
\end{center}

\begin{plainlanguage}
If you check a 4-week experiment every day, roughly 1 in 4 ``significant''
results could be false.  This means you might ship features that have
no real effect --- or even hurt your metrics.
\end{plainlanguage}

\begin{keytakeaway}
\textbf{The peeking problem:} Checking a traditional test repeatedly inflates
the false positive rate far beyond its nominal level.  We need a fundamentally
different kind of test --- one that remains valid no matter when or how often
we look.

That kind of test exists.  It is called a \textbf{sequential test}, and it is
what Eppo (and similar platforms) use under the hood.
\end{keytakeaway}


% ══════════════════════════════════════════════════════════════════════════════
% ACT 2
% ══════════════════════════════════════════════════════════════════════════════
\newpage
\section{Act~2 --- The Eppo Solution\\
         (Schmit \& Miller, 2024)}

\subsection{The Big Idea}

\begin{intuition}
Eppo replaces the standard confidence interval with a \textbf{sequential
confidence interval} --- a band around the estimated treatment effect that
is valid at \emph{every moment} since the experiment started, not just at
one pre-planned analysis time.

The key insight: a sequential CI is wider than a standard CI (you pay a
small price for the freedom to peek), but the error guarantee holds no
matter when you look.  And Eppo reduces the noise in your data using
pre-experiment information, making the CI narrower and experiments faster.
\end{intuition}

\subsection{The Simulation}

\begin{simulation}
\textbf{What is being simulated:}
A single A/B experiment flowing through Eppo's pipeline, stage by stage.

\textbf{Animated pipeline diagram.}  Data flows left-to-right through
labelled stages.  Each stage lights up as it is explained.

\textbf{Stages (rendered as connected boxes):}
\[
  \fbox{Randomise}
  \;\to\;
  \fbox{Collect}
  \;\to\;
  \fbox{Adjust}
  \;\to\;
  \fbox{Estimate}
  \;\to\;
  \fbox{CI}
  \;\to\;
  \fbox{Decide}
\]

\textbf{Step-by-step mode:}
A button labelled \textbf{``Step through pipeline''} pauses at each stage.
At each stage, a distinct visual state is shown:

\begin{center}
\renewcommand{\arraystretch}{1.3}
\begin{tabular}{@{} l p{9cm} @{}}
\toprule
\textbf{Stage} & \textbf{Visual state} \\
\midrule
Randomise & Users appear as dots, randomly coloured blue (control) or
  green (treatment). \\
Collect & Each dot acquires labels: a pre-experiment value $X$ and an
  outcome $Y$.  Tooltip on hover shows both values. \\
Adjust & Dots animate: $Y \to Y^*$.  A regression line overlays each group.
  Dots shift vertically to their residual position.  Spread visibly
  decreases. \\
Estimate & Group means appear as horizontal lines.  The lift displays
  as a labelled gap between lines. \\
CI & A confidence interval band draws around the estimated lift.
  The band narrows as more data arrives. \\
Decide & The CI is colour-coded:
  \textcolor{green!60!black}{\textbf{green}} (entirely above 0),
  \textcolor{red}{\textbf{red}} (entirely below 0), or
  \textcolor{gray}{\textbf{grey}} (crosses 0).  A verdict label appears. \\
\bottomrule
\end{tabular}
\end{center}

\textbf{Side-by-side comparison mode:}
A button labelled \textbf{``Compare with / without adjustment''}
splits the display into two panels using the \emph{same data} ---
making the benefit of noise reduction immediately visible.

\textbf{Controls:}
\begin{itemize}[nosep]
  \item True effect size $\delta$ (0 to 10\%).
  \item Sample size $n$ (100 to 100,000).
  \item Pre-experiment correlation $\rho$ (0 to 0.95): controls how much
        noise the adjustment removes.
  \item Fixed seed toggle (deterministic replay).
  \item ``Run 1,000 experiments'' mode (batch statistics).
\end{itemize}

\textbf{Key takeaway:}
Each pipeline stage transforms the data in a specific, visible way.
Regression adjustment visibly tightens the CI; removing it shows the cost.
\end{simulation}


\subsection{The Pipeline, Step by Step}

\subsubsection*{Step 1: Randomise}

Each user arriving at the website is randomly assigned to \textbf{control}
(existing experience) or \textbf{treatment} (new feature).  Randomisation
ensures the groups are comparable --- any difference in outcomes is caused by
the treatment, not by pre-existing differences.

\subsubsection*{Step 2: Collect data over time}

As users interact with the product, two quantities are recorded for each user:
\begin{itemize}[nosep]
  \item $Y_i$: the \textbf{outcome} during the experiment (e.g.\ purchase
        amount, clicks, session length).
  \item $X_i$: a \textbf{pre-experiment covariate} --- the same metric from
        before the experiment started (e.g.\ purchases in the prior 2 weeks).
\end{itemize}

\subsubsection*{Step 3: Regression adjustment (noise removal)}

Users are wildly different --- some buy \texteuro0, others \texteuro500.
This noise drowns out the treatment effect.  Eppo reduces this noise by
using pre-experiment data.

For each group \textbf{separately}, fit a model predicting $Y$ from $X$:
\[
  \hat{f}_{\text{control}}(X), \qquad \hat{f}_{\text{treatment}}(X)
\]

Then compute the \textbf{adjusted outcome} for each user:
\[
  Y^*_i = Y_i - \hat{f}_g(X_i)
\]
where $g$ is the user's group (control or treatment).

\begin{plainlanguage}
``Predict what each user would have done based on their history, then subtract
that prediction.  What's left is the part you \emph{couldn't} have predicted
--- and that's where the treatment effect hides.''

This is a generalisation of a technique called CUPED (Deng et al., 2013).
The stronger the correlation $\rho$ between pre-experiment and experiment
behaviour, the more noise is removed:
$\Var(Y^*) = \Var(Y)(1 - \rho^2)$.
At $\rho = 0.7$, half the variance is eliminated ---
equivalent to doubling your traffic.
\end{plainlanguage}

\textbf{Safety rule:} The covariate $X$ must come from the
\emph{pre-experiment} period.  Never use data that could be affected by
the treatment.

\subsubsection*{Step 4: Estimate the treatment effect}

Compute the adjusted group means and the \textbf{relative lift}:
\[
  \hat{\tau}(t) = \hat{\mu}_1(t) - \hat{\mu}_0(t), \qquad
  \text{Relative lift}(t) = \frac{\hat{\tau}(t)}{\hat{\mu}_0(t)}
\]

\begin{plainlanguage}
Here $\hat{\mu}_0(t)$ and $\hat{\mu}_1(t)$ are the average adjusted outcomes
in the control and treatment groups at time $t$.  Their difference
$\hat{\tau}(t)$ is the estimated treatment effect.  Dividing by the control
mean gives a percentage: ``the treatment increased purchases by $+3\%$.''

Relative lift is preferred because it is \textbf{scale-invariant} --- a 3\%
lift means the same thing whether the baseline is \texteuro10 or
\texteuro10{,}000.
\end{plainlanguage}

\subsubsection*{Step 5: Estimate the variance}

\[
  \hat{\sigma}_{\hat{\tau}}^2(t)
  = \frac{s_0^2(t)}{n_0(t)} + \frac{s_1^2(t)}{n_1(t)}
\]

\begin{plainlanguage}
Here $s_g^2(t)$ is the sample variance (spread) of the adjusted outcomes in
group~$g$, and $n_g(t)$ is the number of users in that group so far.
The variance is estimated from the data --- not assumed known.
\end{plainlanguage}

\subsubsection*{Step 6: Construct the sequential confidence interval}

This is the key formula.  Instead of a standard CI (which uses the constant
$z_{\alpha/2} \approx 1.96$), Eppo uses a time-varying multiplier:

\[
  \boxed{\text{CI}(t) = \hat{\tau}(t) \;\pm\;
  \hat{\sigma}_{\hat{\tau}}(t) \cdot
  \underbrace{\sqrt{\frac{n + \nu}{n} \cdot
  \log\!\frac{n + \nu}{\nu \alpha^2}}}_{\text{sequential multiplier}}}
\]

\begin{plainlanguage}
\textbf{Unpacking the formula:}
\begin{itemize}[nosep]
  \item $\hat{\tau}(t)$: the estimated treatment effect at time $t$
        (from Step~4).
  \item $\hat{\sigma}_{\hat{\tau}}(t)$: the estimated standard error
        (from Step~5).
  \item $n = n_0(t) + n_1(t)$: total observations so far.
  \item $\alpha$: significance level (e.g.\ 0.05 for a 95\% CI).
  \item $\nu$: a tuning parameter set from the expected total sample size
        $M$: $\nu = M \cdot \hat{\sigma}^2$.  It controls \emph{which}
        sample size the CI is tightest for.
  \item The ``sequential multiplier'' replaces the constant 1.96.  It
        starts higher (e.g.\ $\sim\!3.8$ early on) and slowly decreases
        toward 1.96 --- but never quite reaches it.  This is the
        \textbf{price of peeking}.
\end{itemize}

This formula comes from the \emph{confidence sequence} framework of
Howard et al.\ (2021), specifically the Normal mixture boundary.  The
mathematical guarantee: the true treatment effect lies inside the CI at
\emph{all times simultaneously} with probability $\geq 1 - \alpha$.
\end{plainlanguage}

\subsubsection*{Step 7: Decide}

\begin{center}
\renewcommand{\arraystretch}{1.3}
\begin{tabular}{@{} l l p{6.5cm} @{}}
\toprule
\textbf{What you see} & \textbf{Colour} & \textbf{Action} \\
\midrule
CI entirely above 0
  & \textcolor{green!60!black}{\rule{1em}{1em}} Green
  & Positive effect.  Ship the feature. \\
CI entirely below 0
  & \textcolor{red}{\rule{1em}{1em}} Red
  & Negative effect.  Revert. \\
CI crosses 0, experiment ongoing
  & \textcolor{gray}{\rule{1em}{1em}} Grey
  & Inconclusive.  Keep collecting. \\
CI crosses 0, experiment ended
  & \textcolor{gray}{\rule{1em}{1em}} Grey
  & No significant effect.  Decide pragmatically. \\
\bottomrule
\end{tabular}
\end{center}

\textbf{This decision is valid at whatever time it is made.}  No need to
pre-specify an analysis time.  This solves the peeking problem from Act~1.

\begin{keytakeaway}
\textbf{The Eppo pipeline in one sentence:}
Randomise $\to$ collect outcomes $Y$ and covariates $X$ $\to$
regression-adjust (remove predictable noise) $\to$ estimate lift and variance
$\to$ wrap in a sequential CI $\to$ decide when CI excludes zero.

\textbf{What makes it work:}
\begin{itemize}[nosep]
  \item \textbf{Framework:} Howard et al.\ (2021) confidence sequences.
  \item \textbf{Noise reduction:} Generalised CUPED via per-group regression.
  \item \textbf{Output:} Relative lift with a CI that is valid at any time.
\end{itemize}
\end{keytakeaway}


% ══════════════════════════════════════════════════════════════════════════════
% ACT 3
% ══════════════════════════════════════════════════════════════════════════════
\newpage
\section{Act~3 --- What If You Don't Have Eppo?}

Act~2 described Eppo's sophisticated pipeline.  But many teams run experiments
without a dedicated platform.  This act presents three progressively better
approaches to peek-safe testing that anyone can implement with a spreadsheet
or a few lines of code --- and compares them head-to-head with Eppo.

\subsection{The Simulation}

\begin{simulation}
\textbf{What is being simulated:}
A single A/B experiment monitored at $K$ pre-planned analysis times, showing
four methods side-by-side.

\textbf{Panel~1 --- Four confidence bands:}
The same experiment, with four CI bands evolving over time:
\begin{itemize}[nosep]
  \item \textcolor{red}{\textbf{Red (Bonferroni):}} Widest.  Simple but
        conservative.
  \item \textcolor{orange}{\textbf{Orange (Pocock):}} Constant threshold,
        tighter than Bonferroni.
  \item \textcolor{blue}{\textbf{Blue (O'Brien--Fleming):}} Starts very wide,
        narrows to near-classical at the final analysis.
  \item \textcolor{green!60!black}{\textbf{Green (Eppo / Howard):}}
        Continuous monitoring.  No need to pre-specify $K$.
\end{itemize}

Vertical dashed lines mark the $K$ planned analysis times.  Between analyses,
the group-sequential methods have no valid CI --- only Eppo's band is defined
everywhere.

\textbf{Panel~2 --- 10{,}000 experiments:}
Run 10{,}000 experiments under the null ($\delta = 0$).  Counters show the
false positive rate for each method:
\begin{itemize}[nosep]
  \item Classical CI with $K$ peeks: $\gg 5\%$.
  \item Bonferroni: $\leq 5\%$ (often well below).
  \item Pocock: $\leq 5\%$ (closer to 5\%).
  \item O'Brien--Fleming: $\leq 5\%$.
  \item Eppo: $\leq 5\%$.
\end{itemize}

\textbf{Panel~3 --- Efficiency comparison:}
Under a true effect ($\delta > 0$), show average stopping time for each
method.

\textbf{Controls:}
\begin{itemize}[nosep]
  \item Number of planned peeks $K$ (2 to 20).
  \item True effect size $\delta$ (0 to 10\%).
  \item Significance level $\alpha$.
  \item Fixed seed toggle (deterministic replay).
  \item ``Run 1{,}000 experiments'' mode.
\end{itemize}

\textbf{Key takeaway:}
All four methods control the false positive rate.  The price of simplicity
is wider intervals (slower detection).  Even the simplest method
(Bonferroni) is far better than naive peeking.
\end{simulation}


\subsection{Intuitive Explanation}

\begin{intuition}
You want to check your experiment every Monday.  If you use a standard
95\% CI each time, you'll get false positives far too often (Act~1).

The simplest fix: \textbf{be stricter at each check.}  If you plan to peek
4~times, demand stronger evidence each time.  The three methods below
differ in \emph{how} they distribute that strictness across the peeks.

Think of it as a budget.  You have a total error budget of $\alpha = 5\%$.
\begin{itemize}[nosep]
  \item \textbf{Bonferroni} distributes the budget equally:
        $\alpha/K$ per peek.
  \item \textbf{Pocock} distributes it more cleverly (using the correlation
        between test statistics at different peeks), but still equally across
        peeks.
  \item \textbf{O'Brien--Fleming} front-loads the budget: almost no spending
        early, most of it saved for the final analysis.
\end{itemize}

Eppo's approach (Act~2) doesn't need a fixed budget at all --- it works
for \emph{any} number of peeks, continuously.
\end{intuition}


\subsection{Method 1: Bonferroni Correction (Simplest)}

\subsubsection*{The idea}

Before the experiment starts, commit to checking exactly $K$ times
(e.g.\ every Monday for 4 weeks: $K = 4$).  At each check, use a
significance level of $\alpha/K$ instead of $\alpha$.

\subsubsection*{The recipe}

\begin{enumerate}[nosep]
  \item Choose $K$ (number of planned peeks) and $\alpha$ (e.g.\ 0.05).
  \item At each peek $k = 1, \ldots, K$, compute the standard CI using
        the adjusted significance level:
        \[
          \text{CI}_k = \hat{\tau}(t_k) \;\pm\;
          \hat{\sigma}_{\hat{\tau}}(t_k) \cdot z_{\alpha/(2K)}
        \]
        where $\hat{\tau}(t_k)$ is the estimated treatment effect at peek $k$,
        $\hat{\sigma}_{\hat{\tau}}(t_k)$ is its standard error, and
        $z_{\alpha/(2K)}$ is the critical value from a Normal distribution
        at the adjusted significance level.
  \item If the CI does not cross 0: stop and declare significance.
  \item If $k = K$ and the CI crosses 0: no significant effect.
\end{enumerate}

\begin{center}
\renewcommand{\arraystretch}{1.3}
\begin{tabular}{@{} c c c c @{}}
\toprule
$K$ (peeks) & $\alpha/K$ & $z_{\alpha/(2K)}$ & CI multiplier vs.\ 1.96 \\
\midrule
1  & 0.050  & 1.96 & $1.00\times$ (no correction) \\
2  & 0.025  & 2.24 & $1.14\times$ \\
4  & 0.0125 & 2.50 & $1.28\times$ \\
8  & 0.00625 & 2.73 & $1.39\times$ \\
12 & 0.00417 & 2.87 & $1.46\times$ \\
20 & 0.0025  & 3.02 & $1.54\times$ \\
52 & 0.00096 & 3.29 & $1.68\times$ (weekly for a year) \\
\bottomrule
\end{tabular}
\end{center}

\begin{plainlanguage}
With 4 planned peeks, replace $z = 1.96$ with $z = 2.50$.  Your CI is
28\% wider at each peek, but you can safely check 4 times without
inflating errors.
\end{plainlanguage}

\subsubsection*{Implementation step by step}

Given data arrays \texttt{control[]} and \texttt{treatment[]} at each
scheduled peek:

\begin{enumerate}
  \item \textbf{Compute group statistics:}
  \[
    \bar{Y}_0 = \frac{1}{n_0}\sum_{i=1}^{n_0} Y_{0,i}, \qquad
    s_0^2 = \frac{1}{n_0 - 1}\sum_{i=1}^{n_0}(Y_{0,i} - \bar{Y}_0)^2
  \]
  and analogously for group~1.

  \item \textbf{Compute the estimated treatment effect and its standard
        error:}
  \[
    \hat{\tau} = \bar{Y}_1 - \bar{Y}_0, \qquad
    \text{SE} = \sqrt{\frac{s_0^2}{n_0} + \frac{s_1^2}{n_1}}
  \]

  \item \textbf{Look up the adjusted critical value:}
  \[
    z^* = z_{\alpha/(2K)} = \Phi^{-1}\!\bigl(1 - \tfrac{\alpha}{2K}\bigr)
  \]
  where $\Phi^{-1}$ is the inverse of the standard Normal CDF 
  (e.g.\ \texttt{scipy.stats.norm.ppf} in Python).

  \item \textbf{Construct the CI and decide:}
  \[
    \text{CI} = \bigl[\hat{\tau} - z^* \cdot \text{SE},\;
                       \hat{\tau} + z^* \cdot \text{SE}\bigr]
  \]
  If the CI does not contain 0, stop and declare significance.
  Otherwise, wait until the next scheduled peek.
\end{enumerate}

\subsubsection*{Pros and cons}

\begin{itemize}[nosep]
  \item[\textcolor{green!60!black}{+}] \textbf{Dead simple.}  One line of
        code: change \texttt{1.96} to \texttt{z\_alpha/(2K)}.
  \item[\textcolor{green!60!black}{+}] \textbf{Always valid.}  Works for
        any test statistic, any distribution.
  \item[\textcolor{red}{--}] \textbf{Conservative.}  The CI is wider than
        necessary because Bonferroni ignores the correlation between test
        statistics at successive peeks.
  \item[\textcolor{red}{--}] \textbf{Scales poorly.}  As $K$ grows, the
        correction becomes increasingly harsh.
\end{itemize}


\subsection{Method 2: Pocock Boundaries (1977)}

\subsubsection*{The idea}

Like Bonferroni, pre-specify $K$ equally-spaced analysis times.  But instead
of Bonferroni's crude $\alpha/K$ split, use a constant critical value $c_P$
that accounts for the correlation between test statistics at different peeks.

Because the test statistics $Z_1, Z_2, \ldots, Z_K$ at successive analyses
share overlapping data, they are positively correlated.  Pocock's method
exploits this to derive a tighter threshold.

\subsubsection*{The recipe}

\begin{enumerate}[nosep]
  \item Choose $K$ and $\alpha$.
  \item Look up (or compute) the Pocock critical value $c_P(K, \alpha)$.
  \item At each peek $k$, compute the $z$-statistic:
        $Z_k = \hat{\tau}(t_k) / \hat{\sigma}_{\hat{\tau}}(t_k)$.
  \item If $|Z_k| > c_P$: stop and declare significance.
\end{enumerate}

\begin{center}
\renewcommand{\arraystretch}{1.3}
\begin{tabular}{@{} c c c c @{}}
\toprule
$K$ & $c_P$ (Pocock) & $z_{\alpha/(2K)}$ (Bonf.) & Improvement \\
\midrule
2  & 2.18 & 2.24 & 3\% tighter \\
4  & 2.36 & 2.50 & 6\% tighter \\
8  & 2.51 & 2.73 & 8\% tighter \\
12 & 2.60 & 2.87 & 9\% tighter \\
20 & 2.69 & 3.02 & 11\% tighter \\
\bottomrule
\end{tabular}
\end{center}

\begin{plainlanguage}
Pocock gives a constant threshold that is a few percent lower than
Bonferroni's.  The key advantage: the threshold is the same at every peek,
so it is almost as simple to use as Bonferroni.
\end{plainlanguage}

\subsubsection*{Implementation step by step}

The Pocock critical value $c_P$ is found by solving the following equation
numerically.  Let $\mathbf{Z} = (Z_1, \ldots, Z_K)$ be the vector of
$z$-statistics at the $K$ equally-spaced analyses.  Under the null, these
follow a multivariate Normal distribution with correlation matrix:
\[
  \Cor(Z_j, Z_k) = \sqrt{\frac{\min(j,k)}{\max(j,k)}}
  \qquad \text{for } 1 \leq j, k \leq K
\]

The Pocock boundary $c_P$ is the unique value satisfying:
\[
  \Prob\!\bigl(\max_{1 \leq k \leq K} |Z_k| > c_P\bigr) = \alpha
\]

This is computed via multivariate Normal integration (see Appendix~A for
Python code using \texttt{scipy.stats.multivariate\_normal}).

\textbf{Once $c_P$ is known}, the decision rule at each peek is identical
to Bonferroni, but with $c_P$ replacing $z_{\alpha/(2K)}$:

\begin{enumerate}
  \item At peek $k$, compute:
  \[
    Z_k = \frac{\hat{\tau}(t_k)}{\text{SE}(t_k)}
  \]
  \item If $|Z_k| > c_P$: reject the null.  Otherwise: continue.
\end{enumerate}


\subsection{Method 3: O'Brien--Fleming Boundaries (1979)}

\subsubsection*{The idea}

Instead of a constant threshold, use a threshold that \emph{decreases} over
time.  At the first peek, demand very strong evidence; at the last peek,
demand nearly the same evidence as a classical test.

\subsubsection*{The recipe}

At peek $k$ (of $K$ total), use the critical value:
\[
  c_k^{\text{OBF}} \;\approx\; z_{\alpha/2} \cdot \sqrt{K/k}
\]
where $z_{\alpha/2} \approx 1.96$ is the standard Normal critical value.

\begin{center}
\renewcommand{\arraystretch}{1.3}
\begin{tabular}{@{} c c c c c @{}}
\toprule
Peek $k$ (of $K{=}4$) & $c_k^{\text{OBF}}$ & $c_P$ (Pocock) &
  $z_{\alpha/(2K)}$ (Bonf.) & $z_{\alpha/2}$ (classical) \\
\midrule
1 (25\% of data) & 4.05 & 2.36 & 2.50 & 1.96 \\
2 (50\% of data) & 2.86 & 2.36 & 2.50 & 1.96 \\
3 (75\% of data) & 2.34 & 2.36 & 2.50 & 1.96 \\
4 (100\% of data) & 2.02 & 2.36 & 2.50 & 1.96 \\
\bottomrule
\end{tabular}
\end{center}

\begin{plainlanguage}
O'Brien--Fleming is extremely conservative early (threshold 4.05 at the
first peek --- you'd almost never stop that early) but barely penalises
the final analysis (2.02 vs.\ 1.96 --- only 3\% wider).

This is the most efficient approach: if you run the experiment to completion,
you pay almost nothing for the option to peek.
\end{plainlanguage}

\subsubsection*{Implementation step by step}

The O'Brien--Fleming boundary is the simplest to compute of the three
``non-trivial'' methods --- it has a closed-form approximation.

\begin{enumerate}
  \item \textbf{Before the experiment:} Fix $K$ and $\alpha$.  At each peek
        $k$, the fraction of total data collected is
        $t_k = k / K$ (assuming equal spacing).

  \item \textbf{At peek $k$, compute the critical value:}
  \[
    c_k^{\text{OBF}} = \frac{z_{\alpha/2}}{\sqrt{t_k}}
    = z_{\alpha/2} \cdot \sqrt{\frac{K}{k}}
  \]
  This is the Lan--DeMets spending function approximation to the original
  O'Brien--Fleming boundary.

  \item \textbf{Compute the $z$-statistic} (same as Bonferroni and Pocock):
  \[
    Z_k = \frac{\hat{\tau}(t_k)}{\text{SE}(t_k)}
  \]

  \item \textbf{Decide:}
  If $|Z_k| > c_k^{\text{OBF}}$: reject.  Otherwise: continue to peek $k+1$.

  \item \textbf{Convert to a confidence interval} (optional but recommended
        for interpretation):
  \[
    \text{CI}_k = \hat{\tau}(t_k) \;\pm\; c_k^{\text{OBF}} \cdot
    \text{SE}(t_k)
  \]
  Note: the CI width \emph{changes} across peeks because $c_k^{\text{OBF}}$
  varies.
\end{enumerate}

\begin{plainlanguage}
\textbf{Practical tip:}
The exact O'Brien--Fleming boundaries (computed via multivariate Normal
integration, like Pocock) differ slightly from the $z_{\alpha/2}/\sqrt{t_k}$
approximation.  The approximation is accurate for $K \leq 10$ and is what
most software implements.  For larger $K$, use a library
(e.g.\ \texttt{statsmodels.stats.multitest} or the R package \texttt{gsDesign}).
\end{plainlanguage}


\subsection{Head-to-Head Comparison}

\begin{center}
\renewcommand{\arraystretch}{1.3}
\begin{tabular}{@{} l c c c c @{}}
\toprule
& \textbf{Bonferroni} & \textbf{Pocock} & \textbf{OBF} & \textbf{Eppo} \\
\midrule
Pre-specify $K$?
  & Yes & Yes & Yes & No \\
Threshold type
  & Constant & Constant & Decreasing & Continuous \\
CI at peek 1 ($K{=}4$)
  & $2.50\times\hat{\sigma}$ & $2.36\times\hat{\sigma}$
  & $4.05\times\hat{\sigma}$ & $\sim\!3.8\times\hat{\sigma}$ \\
CI at final analysis ($K{=}4$)
  & $2.50\times\hat{\sigma}$ & $2.36\times\hat{\sigma}$
  & $2.02\times\hat{\sigma}$ & $\sim\!2.3\times\hat{\sigma}$ \\
Valid between peeks?
  & No & No & No & Yes \\
Variance reduction?
  & Manual & Manual & Manual & Built-in \\
Implementation
  & 1 line & Lookup table & Formula & Platform \\
\bottomrule
\end{tabular}
\end{center}

\begin{keytakeaway}
\textbf{Which method should you use?}
\begin{itemize}[nosep]
  \item \textbf{No platform, minimal effort:} Bonferroni.  Commit to $K$
        peeks.  Replace 1.96 with $z_{\alpha/(2K)}$.  Done.
  \item \textbf{No platform, moderate effort:} O'Brien--Fleming.  Nearly as
        efficient as Eppo at the final analysis, but requires pre-specifying
        $K$ and cannot peek between scheduled analyses.
  \item \textbf{Need continuous monitoring or variance reduction:} Use a
        platform like Eppo.  Continuous validity, regression adjustment, and
        no need to pre-commit to $K$.
\end{itemize}

\textbf{The bottom line:} even the crudest correction (Bonferroni) is
\emph{vastly} better than naive peeking.  If you are checking results
without any correction, you should start with Bonferroni \emph{today}.
\end{keytakeaway}


% ══════════════════════════════════════════════════════════════════════════════
\newpage
\section{Summary}

\begin{center}
\renewcommand{\arraystretch}{1.3}
\begin{tabular}{@{} c l p{8cm} @{}}
\toprule
\textbf{Act} & \textbf{Topic} & \textbf{Key message} \\
\midrule
1 & The Peeking Problem
  & Checking a standard test repeatedly inflates false positives to
    20--30\%.  You need a different kind of test. \\
2 & The Eppo Solution
  & A sequential CI is valid at every moment.  Regression adjustment
    removes noise, making experiments faster. \\
3 & DIY Alternatives
  & Bonferroni (simplest), Pocock, O'Brien--Fleming --- three ways to
    peek safely without a platform. \\
\bottomrule
\end{tabular}
\end{center}

\bigskip
\textbf{The core trade-off:} if you want to peek, you must pay --- your
confidence interval must be wider.  The question is how much wider, and the
answer depends on how sophisticated a method you use:

\begin{center}
\begin{tabular}{@{} l l @{}}
\toprule
\textbf{Method} & \textbf{Price of peeking} \\
\midrule
No correction (naive) & Wrong answers (20--30\% false positives) \\
Bonferroni & CI $\sim\!28$--$68\%$ wider (depending on $K$) \\
Pocock & CI $\sim\!20$--$37\%$ wider \\
O'Brien--Fleming & CI $\sim\!3\%$ wider at final analysis \\
Eppo (sequential CI) & CI $\sim\!10$--$40\%$ wider, valid \emph{continuously} \\
\bottomrule
\end{tabular}
\end{center}


% ══════════════════════════════════════════════════════════════════════════════
\newpage
\section{References}

\begin{enumerate}
  \item \textbf{Deng, A., Xu, Y., Kohavi, R., \& Walker, T.} (2013).
        Improving the sensitivity of online controlled experiments by
        utilizing pre-experiment data.
        \textit{Proc.\ WSDM}, 123--132.
        \url{https://doi.org/10.1145/2433396.2433413}
        --- The original CUPED paper.  Variance reduction via pre-experiment
        covariates.

  \item \textbf{Howard, S.\,R., Ramdas, A., McAuliffe, J., \& Sekhon, J.}
        (2021).
        Time-uniform, nonparametric, nonasymptotic confidence sequences.
        \textit{Annals of Statistics}, 49(2), 1055--1080.
        \url{https://doi.org/10.1214/20-AOS1991}
        --- The confidence sequence framework that Eppo adopted.

  \item \textbf{Pocock, S.\,J.} (1977).
        Group sequential methods in the design and analysis of clinical
        trials.
        \textit{Biometrika}, 64(2), 191--199.

  \item \textbf{O'Brien, P.\,C.\ \& Fleming, T.\,R.} (1979).
        A multiple testing procedure for clinical trials.
        \textit{Biometrics}, 35(3), 549--556.

  \item \textbf{Schmit, M.\ \& Miller, T.} (2024).
        Sequential confidence intervals for comparing two proportions with
        variance reduction.
        \textit{Eppo Technical Report}.
        --- Eppo's implementation combining Howard et al.'s confidence
        sequences with generalised CUPED.
\end{enumerate}


% ══════════════════════════════════════════════════════════════════════════════
% APPENDIX A
% ══════════════════════════════════════════════════════════════════════════════
\newpage
\appendix
\section{Python Code}

This appendix provides self-contained Python code for:
\begin{enumerate}[nosep]
  \item A simulation demonstrating the peeking problem (Act~1).
  \item Implementation snippets for each DIY method (Act~3).
\end{enumerate}

All code requires \texttt{numpy} and \texttt{scipy}.  Install with:
\texttt{pip install numpy scipy}.


% ──────────────────────────────────────────────────────────────────────────────
\subsection{Simulation: The Peeking Problem}

This simulation runs many A/B experiments under the null (no true effect)
and measures how often naive peeking produces a false positive, versus
checking only at the end.

\begin{lstlisting}
import numpy as np
from scipy import stats

def simulate_peeking(
    n_experiments: int = 10_000,
    n_obs_per_group: int = 500,
    peek_interval: int = 10,
    alpha: float = 0.05,
    seed: int = 42,
) -> dict:
    """Simulate the peeking problem under the null hypothesis.

    Returns a dict with false positive rates for:
      - 'naive_peeking': check every peek_interval observations
      - 'single_check':  check only at the end
    """
    rng = np.random.default_rng(seed)
    peek_times = list(range(peek_interval,
                            n_obs_per_group + 1,
                            peek_interval))

    false_pos_peek = 0
    false_pos_end = 0

    for _ in range(n_experiments):
        # Generate data from identical distributions (null)
        control = rng.normal(0, 1, n_obs_per_group)
        treatment = rng.normal(0, 1, n_obs_per_group)

        # --- Naive peeking ---
        ever_significant = False
        for t in peek_times:
            ctrl_t = control[:t]
            trt_t = treatment[:t]
            _, p_value = stats.ttest_ind(ctrl_t, trt_t)
            if p_value < alpha:
                ever_significant = True
                break  # stop at first significance
        if ever_significant:
            false_pos_peek += 1

        # --- Single check at the end ---
        _, p_value_end = stats.ttest_ind(control, treatment)
        if p_value_end < alpha:
            false_pos_end += 1

    return {
        "naive_peeking": false_pos_peek / n_experiments,
        "single_check": false_pos_end / n_experiments,
        "n_experiments": n_experiments,
        "n_peeks": len(peek_times),
    }

# --- Run it ---
results = simulate_peeking()
print(f"False positive rate (single check): "
      f"{results['single_check']:.1%}")
print(f"False positive rate (peeking every "
      f"{10} obs, {results['n_peeks']} peeks): "
      f"{results['naive_peeking']:.1%}")
\end{lstlisting}

\begin{plainlanguage}
Expected output (approximately):
\begin{verbatim}
False positive rate (single check): 5.0%
False positive rate (peeking every 10 obs, 50 peeks): 27.3%
\end{verbatim}
\end{plainlanguage}


% ──────────────────────────────────────────────────────────────────────────────
\subsection{Implementation: Bonferroni Correction}

\begin{lstlisting}
import numpy as np
from scipy import stats

def bonferroni_test(
    control: np.ndarray,
    treatment: np.ndarray,
    K: int,
    alpha: float = 0.05,
) -> dict:
    """Run a single-peek Bonferroni-corrected test.

    Call this function at each of the K scheduled peeks
    with the data collected so far.

    Returns:
      - 'reject': whether to reject the null
      - 'ci': the confidence interval for tau_hat
      - 'tau_hat': estimated treatment effect
      - 'z_star': adjusted critical value
    """
    n0, n1 = len(control), len(treatment)
    tau_hat = treatment.mean() - control.mean()
    se = np.sqrt(control.var(ddof=1) / n0
                 + treatment.var(ddof=1) / n1)

    # Bonferroni-adjusted critical value
    z_star = stats.norm.ppf(1 - alpha / (2 * K))

    ci_lower = tau_hat - z_star * se
    ci_upper = tau_hat + z_star * se
    reject = ci_lower > 0 or ci_upper < 0

    return {
        "reject": reject,
        "ci": (ci_lower, ci_upper),
        "tau_hat": tau_hat,
        "z_star": z_star,
        "se": se,
    }

# --- Example usage ---
rng = np.random.default_rng(0)
K = 4  # plan to peek 4 times
N_total = 2000  # total obs per group
peek_every = N_total // K

control_all = rng.normal(10, 3, N_total)
treatment_all = rng.normal(10.3, 3, N_total)  # small effect

for k in range(1, K + 1):
    n = k * peek_every
    result = bonferroni_test(
        control_all[:n], treatment_all[:n], K=K
    )
    ci = result["ci"]
    status = "REJECT" if result["reject"] else "continue"
    print(f"Peek {k}/{K} (n={n}): "
          f"tau={result['tau_hat']:.3f}, "
          f"CI=[{ci[0]:.3f}, {ci[1]:.3f}], "
          f"z*={result['z_star']:.3f} -> {status}")
\end{lstlisting}


% ──────────────────────────────────────────────────────────────────────────────
\subsection{Implementation: Pocock Boundaries}

\begin{lstlisting}
import numpy as np
from scipy import stats, optimize

def pocock_critical_value(
    K: int, alpha: float = 0.05, n_sim: int = 200_000,
    seed: int = 42,
) -> float:
    """Find the Pocock boundary c_P by simulation.

    Simulates the multivariate Normal distribution of
    (Z_1, ..., Z_K) under the null and finds c_P such that
    P(max |Z_k| > c_P) = alpha.

    For production use, replace with exact multivariate
    Normal integration (e.g. mvnorm from R or PyMVPA).
    """
    rng = np.random.default_rng(seed)

    # Build the correlation matrix: Cor(Z_j, Z_k) = sqrt(min/max)
    fracs = np.arange(1, K + 1) / K
    corr = np.zeros((K, K))
    for j in range(K):
        for k in range(K):
            corr[j, k] = np.sqrt(
                min(fracs[j], fracs[k])
                / max(fracs[j], fracs[k])
            )

    # Simulate Z vectors
    samples = rng.multivariate_normal(
        np.zeros(K), corr, size=n_sim
    )
    max_abs_z = np.abs(samples).max(axis=1)

    # Find c_P: the (1 - alpha) quantile of max|Z|
    c_P = np.quantile(max_abs_z, 1 - alpha)
    return float(c_P)

def pocock_test(
    control: np.ndarray,
    treatment: np.ndarray,
    c_P: float,
) -> dict:
    """Test at a single peek using the Pocock boundary."""
    n0, n1 = len(control), len(treatment)
    tau_hat = treatment.mean() - control.mean()
    se = np.sqrt(control.var(ddof=1) / n0
                 + treatment.var(ddof=1) / n1)
    z_stat = tau_hat / se
    reject = abs(z_stat) > c_P

    ci_lower = tau_hat - c_P * se
    ci_upper = tau_hat + c_P * se

    return {
        "reject": reject,
        "z_stat": z_stat,
        "ci": (ci_lower, ci_upper),
        "tau_hat": tau_hat,
    }

# --- Example usage ---
K = 4
c_P = pocock_critical_value(K)
print(f"Pocock critical value for K={K}: {c_P:.3f}")

rng = np.random.default_rng(0)
N_total = 2000
peek_every = N_total // K

control_all = rng.normal(10, 3, N_total)
treatment_all = rng.normal(10.3, 3, N_total)

for k in range(1, K + 1):
    n = k * peek_every
    result = pocock_test(
        control_all[:n], treatment_all[:n], c_P=c_P
    )
    ci = result["ci"]
    status = "REJECT" if result["reject"] else "continue"
    print(f"Peek {k}/{K} (n={n}): "
          f"Z={result['z_stat']:.3f}, "
          f"CI=[{ci[0]:.3f}, {ci[1]:.3f}] -> {status}")
\end{lstlisting}


% ──────────────────────────────────────────────────────────────────────────────
\subsection{Implementation: O'Brien--Fleming Boundaries}

\begin{lstlisting}
import numpy as np
from scipy import stats

def obf_critical_value(k: int, K: int,
                       alpha: float = 0.05) -> float:
    """O'Brien-Fleming critical value at peek k of K.

    Uses the Lan-DeMets spending function approximation:
        c_k = z_{alpha/2} / sqrt(k / K)
    """
    z_alpha2 = stats.norm.ppf(1 - alpha / 2)
    t_k = k / K  # information fraction
    return z_alpha2 / np.sqrt(t_k)

def obf_test(
    control: np.ndarray,
    treatment: np.ndarray,
    k: int,
    K: int,
    alpha: float = 0.05,
) -> dict:
    """Test at peek k of K using O'Brien-Fleming."""
    n0, n1 = len(control), len(treatment)
    tau_hat = treatment.mean() - control.mean()
    se = np.sqrt(control.var(ddof=1) / n0
                 + treatment.var(ddof=1) / n1)
    z_stat = tau_hat / se
    c_k = obf_critical_value(k, K, alpha)
    reject = abs(z_stat) > c_k

    ci_lower = tau_hat - c_k * se
    ci_upper = tau_hat + c_k * se

    return {
        "reject": reject,
        "z_stat": z_stat,
        "c_k": c_k,
        "ci": (ci_lower, ci_upper),
        "tau_hat": tau_hat,
    }

# --- Example usage ---
K = 4
rng = np.random.default_rng(0)
N_total = 2000
peek_every = N_total // K

control_all = rng.normal(10, 3, N_total)
treatment_all = rng.normal(10.3, 3, N_total)

for k in range(1, K + 1):
    n = k * peek_every
    result = obf_test(
        control_all[:n], treatment_all[:n], k=k, K=K
    )
    ci = result["ci"]
    status = "REJECT" if result["reject"] else "continue"
    print(f"Peek {k}/{K} (n={n}): "
          f"Z={result['z_stat']:.3f}, "
          f"c_k={result['c_k']:.3f}, "
          f"CI=[{ci[0]:.3f}, {ci[1]:.3f}] -> {status}")
\end{lstlisting}


% ──────────────────────────────────────────────────────────────────────────────
\subsection{Simulation: Comparing All Methods}

This simulation runs all three methods (plus naive peeking) on the same
experiments under the null, and reports their false positive rates.

\begin{lstlisting}
import numpy as np
from scipy import stats

def compare_methods(
    n_experiments: int = 10_000,
    n_obs_per_group: int = 2000,
    K: int = 4,
    alpha: float = 0.05,
    seed: int = 42,
) -> dict:
    """Compare false positive rates of all methods."""
    rng = np.random.default_rng(seed)
    peek_every = n_obs_per_group // K

    # Pre-compute critical values
    z_bonf = stats.norm.ppf(1 - alpha / (2 * K))
    c_P = pocock_critical_value(K, alpha)

    counts = {
        "naive": 0,
        "bonferroni": 0,
        "pocock": 0,
        "obf": 0,
    }

    for _ in range(n_experiments):
        ctrl = rng.normal(0, 1, n_obs_per_group)
        trt = rng.normal(0, 1, n_obs_per_group)

        rej = {m: False for m in counts}

        for k in range(1, K + 1):
            n = k * peek_every
            c, t = ctrl[:n], trt[:n]
            tau = t.mean() - c.mean()
            se = np.sqrt(c.var(ddof=1)/n + t.var(ddof=1)/n)
            z = abs(tau / se)

            # Naive: standard z > 1.96
            if z > stats.norm.ppf(1 - alpha / 2):
                rej["naive"] = True

            # Bonferroni
            if z > z_bonf:
                rej["bonferroni"] = True

            # Pocock
            if z > c_P:
                rej["pocock"] = True

            # O'Brien-Fleming
            c_k = stats.norm.ppf(1 - alpha/2) / np.sqrt(k/K)
            if z > c_k:
                rej["obf"] = True

        for m, rejected in rej.items():
            if rejected:
                counts[m] += 1

    return {m: c / n_experiments for m, c in counts.items()}

# --- Run it ---
rates = compare_methods(K=4)
print(f"{'Method':<15} {'FPR':>8}")
print("-" * 25)
for method, fpr in rates.items():
    label = method.upper().replace("_", " ")
    marker = " ***" if fpr > 0.06 else ""
    print(f"{label:<15} {fpr:>7.1%}{marker}")
\end{lstlisting}

\begin{plainlanguage}
Expected output (approximately):
\begin{verbatim}
Method              FPR
-------------------------
NAIVE            14.2% ***
BONFERRONI        3.1%
POCOCK            4.5%
OBF               4.8%
\end{verbatim}
All corrected methods stay at or below 5\%.  Naive peeking exceeds it by
a wide margin, even with only $K = 4$ peeks.
\end{plainlanguage}


\end{document}
